% AI-written, AI cross-reviewed; no human sign-off; not a claim that RH is proved
% RH Notes blueprint source excerpts for N188
% Verbatim excerpts in the notation of the parent blueprint. The labelled editions and source ranges appear below.
% See the unit YAML for the current proof status. Earlier-round passages retain their original wording.

% Start of atlas edition annotation. The following block is outside the original excerpt.
\paragraph{Status of named inputs in this reading.}
\noindent\nolinkurl{GlobalTheory.Published.SiegelOneVector.IsSchwartzBruhat} --- hypothesis at this head (\path{Development/AbelianVarietyTheory/PublishedTheorems/GlobalTheta/IsSchwartzBruhatOfProd.lean:22--26}); not compared.\par
\noindent\nolinkurl{GlobalTheory.T.siegel_weil_measure_of_rank_small_or_zero} --- in-repo constant (\path{CompanionTheory/SpectralGeometry/LocalGlobalProofs/AdelicAnalysis/SiegelWeilMeasureOfRankSmallOrZeroStatement.lean:7--8}); not compared.\par
\noindent\nolinkurl{GlobalTheory.T.siegel_weil_hermitian_coefficient} --- in-repo constant (\path{Development/WeilBound/RepresentationTheoryGL/RegularisedSiegelWeil/ThetaIdempotents/ThetaCoefficients/SiegelWeilHermitianCoefficientStatement.lean:9--10}); not compared.\par
\noindent\nolinkurl{GlobalTheory.T.archimedean_siegel_weil_integral_identity} --- in-repo constant (\path{Development/AutomorphicSpectrum/Proofs/BlockDecomposition/ArchimedeanSiegelWeilIntegralIdentityStatement.lean:15--22}); not compared.\par
\noindent\nolinkurl{GlobalTheory.T.finite_adelic_siegel_weil_integral_identity} --- in-repo constant (\path{Development/AutomorphicSpectrum/Proofs/BlockDecomposition/FiniteAdelicSiegelWeilIntegralIdentityStatement.lean:15--22}); not compared.\par
External theorem numbers in the original excerpt are not verified here unless a separate edition-specific source record documents the comparison.
% End of atlas edition annotation.

\section*{N188. Nonsingular Hermitian Gram coefficient}

% Parent source SHA-256: 44e3b079775c8d7cc0d43b0314bf43a3959be56b2bbafa7f26e5441a1a9163de
% blueprint/ch_T.tex lines 481 through 523; r5:lit:N188
\paragraph{The nonsingular coefficient and its normalization (N187--N189).}\label{r5:lit:N188}
Let $E/F$ be the given CM extension, $V=(E^M,J)$ nondegenerate, and
$Q(X)$ the Hermitian Gram matrix of $n$ vectors. Write
$[b,T]=\tfrac12\operatorname{Tr}_{E/F}\operatorname{tr}(bT)$.
The two functions used in N188 and N189 are
\[
 F^*_\Phi(b)=\int\Phi(X)\psi([b,Q(X)])\,dX,\qquad
 F_\Phi(T)=\int F^*_\Phi(b)\psi(-[b,T])\,db .
\]
The source fixes a nontrivial character of $\mathbb A_F/F$ and self-dual
measures for the indicated trace pairings (\path{IsCharAF}, \path{SdW},
\path{SdHerm}, \path{SWXSDefs.lean:94--120,160--165,213--214,292--299};
\path{starN}, \path{coefN}, \path{ThetaPartD.lean:17--27}).
These are Weil's invariant map and its dual variable for the adjoint
involution on $\operatorname{End}_E(V)$ acting on $V^n$. The involution
has type $\mathrm I_2$ and $\epsilon=1/2$, so his condition
$M>2n+4\epsilon-2$ is exactly $M>2n$. On Hermitian matrices
$\operatorname{tr}(bT)\in\mathbb A_F$, which explains the half-trace in the
coordinate formula. Rational changes of bases or a rational scalar in
the rational pairing change the product of local measure factors by
one, by the product formula. The global self-dual measure has rational
quotient volume one (Weil, \S3). The Schwartz functions in his \S41
include the source's finite sums of archimedean Schwartz functions times
locally constant, compactly supported finite parts
(\path{IsSchwartzBruhatOfProd.lean:18--27}).

Weil proves $F^*_\Phi\in L^1$ and identifies $F_\Phi(T)$ with integration
against a measure on the Gram fibre. His individual-fibre identity for
invariant quotient probability then identifies it with the average of
$\sum_{Q(x)=T}\Phi(h^{-1}x)$ over $U(V)(F)\backslash U(V)(\mathbb A_F)$.
When $\det T\ne0$ every tuple in this sum has full rank. The inverse
changes Weil's right quotient to the left quotient in the source, and
the factor $\mu(\mathcal F)^{-1}$ normalizes its Haar fundamental domain
to probability (\path{thetaMeas}, \path{IsEulerProd.lean:84--87}).
For $n=0$ both integrations reduce to evaluation at the unique tuple.
This gives N189 and N188, respectively, with N187 an alias of N188
(\path{integrable_starN}, \path{SiegelWeilConvergentRange.lean:20--82};
\path{coefN_eq_theta}, \path{SiegelWeilHermitianFormula.lean:303--336};
\path{siegel_weil_hermitian_formula}, \path{SiegelWeilHermitianFormulaStatement.lean:12--22}).
See \cite[\S\S3,41,52, Theorem~5]{T-Weil1965}. Yamana's coefficient
version also has constant one since the Witt index $r$ is at most $M/2$, so
$M-r>n$. His archimedean test space in this statement consists of
Fock polynomials \cite[pp.~236--237, Lemma~3.1]{T-Yamana2013}.

